% RESULTADO_RECUPERADO from II REV09. See MAPA_PROCEDENCIA.json.
% Selected lines and explicit mechanical changes; no new proof.
\section{Area operator and reconstruction of sectorial information}
\label{v:ant:sec:area}
Regrouping six trits constructs a bijection between
$(\Z/9\Z)^3$ and $(\Z/27\Z)^2$, both of cardinality $729$.
The bijection is not presented as an additive isomorphism. On a
$9\times9$ boundary block, the source decomposes 36 links into two
sets of 18, labeled by channels $90$ and $120$.

\subsection{Tritic regrouping and boundary incidence}
\label{v:ant:sec:ii-area-incidencia-explicita}

The correspondence preserves an ordered word, in addition to its cardinality.
Write each coordinate \(x_j\in\mathbb Z/9\mathbb Z\), for
\(j=1,2,3\), through its unique positional representatives
\(x_j=r_{2j-1}+3r_{2j}\), with \(r_k\in\{0,1,2\}\). The map is
\begin{equation}
 \beta_{729}(x_1,x_2,x_3)
 =\bigl(r_1+3r_2+9r_3,\ r_4+3r_5+9r_6\bigr)
 \in(\mathbb Z/27\mathbb Z)^2.
 \label{v:ant:eq:ii-area-beta729}
\end{equation}

\begin{proposition}[Positional correspondence between interior and boundary]
\label{v:ant:prop:ii-area-beta729}
The map \(\beta_{729}\) is bijective and preserves the order of the
six trits. Its inverse extracts the two triples of digits and regroups them
into three pairs.
\end{proposition}
\begin{proof}
The positional expansion of the representatives \(0,\ldots,8\) and
\(0,\ldots,26\) is unique. Extracting the six digits, regrouping them, and
extracting them again returns the same word. The inverse operation
thus restores each of the three initial coordinates.
\end{proof}

Carry preserves its position in the enriched word. The additive
operations of the two packagings have different boundaries: for example,
\(\beta_{729}(0,2,0)=(18,0)\) and
\(\beta_{729}(0,1,0)=(9,0)\), whose sum in the order-twenty-seven
torus is \((0,0)\), whereas
\(\beta_{729}(0,3,0)=(0,1)\). This distinction allows the positional
correspondence to be reused while keeping the operations typed.

In the torus \((\mathbb Z/27\mathbb Z)^2\), the block
\(B_\partial=\{0,\ldots,8\}^2\) has the oriented link sets
\begin{align}
 \partial B_{90}
 &=\{((-1,j),(0,j)),\ ((8,j),(9,j)):0\le j\le8\},\\
 \partial B_{120}
 &=\{((i,-1),(i,0)),\ ((i,8),(i,9)):0\le i\le8\}.
 \label{v:ant:eq:ii-area-enlaces}
\end{align}
All coordinates are read modulo twenty-seven. Each collection contains
eighteen links, and the two are disjoint. Their labels preserve the
incidence channels; the number of links is \(18+18\).

\begin{theorem}[Minimum capacity of the cubic cut]
\label{v:ant:thm:ii-area-corte-minimo}
In the cubic graph with vertices \(\{0,\ldots,N-1\}^3\), with
\(N\ge2\), links between neighbors, and unit capacity per link,
the minimum cut between two opposite faces has capacity \(N^2\).
For \(N=9\), a maximally mixed product state on the
trits of a minimum layer has information
\begin{equation}
 \operatorname{cap}_{\min}=81,\qquad s_{\rm corte}=81\log3.
 \label{v:ant:eq:ii-area-corte-informacion}
\end{equation}
\end{theorem}
\begin{proof}
The \(N^2\) lines parallel to the normal of the faces form edge-disjoint
paths. Every separating cut intercepts each path, so
it contains at least \(N^2\) links. The links crossing
a transverse layer constitute a cut of that cardinality. Each uniform trit
contributes \(\log3\) nats; entropy is additive for the product
state of the eighty-one trits of the layer.
\end{proof}

The capacity of the cubic cut and the state of the thirty-six links
in \eqref{v:ant:eq:ii-area-enlaces} belong to observables defined on
different cuts. The genealogy preserves both domains and their measures.

\subsection{Transport of the collective incidence modes}
\label{v:ant:sec:ii-area-transporte-colectivo}

The spaces \(\mathcal F_6\) and \(\mathcal F_8\) of
Section~\ref{v:ant:sec:barbero} provide the uniform vectors
\[
 u_{90}=\frac1{\sqrt{90}}\sum_{f\in\mathcal F_6}(\delta_f,0),
 \qquad
 u_{120}=\frac1{\sqrt{120}}\sum_{f\in\mathcal F_8}(0,\delta_f)
\]
in \(\ell^2(\mathcal F_6)\oplus\ell^2(\mathcal F_8)\).
In \(\ell^2(\partial B_{90})\oplus\ell^2(\partial B_{120})\),
construct
\[
 v_{90}=\frac1{\sqrt{18}}\sum_{e\in\partial B_{90}}(\delta_e,0),
 \qquad
 v_{120}=\frac1{\sqrt{18}}\sum_{e\in\partial B_{120}}(0,\delta_e).
\]
The two pairs are orthonormal. Denote their respective spanned planes
by \(\mathscr H_{\rm inc}^{\rm coll}\) and
\(\mathscr H_\partial^{\rm coll}\).

\begin{theorem}[Collective intertwining of the sectorial labels]
\label{v:ant:thm:ii-area-entrelazamiento-colectivo}
There is a unique surjective linear isometry between these planes such that
\begin{equation}
 \mathcal J_{6\to8}u_m=v_m\quad(m=90,120).
 \label{v:ant:eq:ii-area-j-colectivo}
\end{equation}
For the operators
\[
 \begin{aligned}
 D_{\rm inc}&=90|u_{90}\rangle\langle u_{90}|
             +120|u_{120}\rangle\langle u_{120}|,\\
 D_\partial&=90|v_{90}\rangle\langle v_{90}|
             +120|v_{120}\rangle\langle v_{120}|,
 \end{aligned}
\]
the following identity holds:
\begin{equation}
 \mathcal J_{6\to8}D_{\rm inc}
 =D_\partial\mathcal J_{6\to8}.
 \label{v:ant:eq:ii-area-j-entrelaza}
\end{equation}
The sectorial projectors and their oriented combination with
coefficients \(12:-1\) are also transported.
\end{theorem}
\begin{proof}
The image of a basis fixes the unique linear map. Norms and
inner products are preserved because both bases are orthonormal.
Identity \eqref{v:ant:eq:ii-area-j-entrelaza} holds on each vector
\(u_m\), since its two sides are \(m v_m\). The projectors are
expressed as \((120I-D)/30\) and \((D-90I)/30\); their linear functional
calculus and any combination of the two preserve intertwining.
\end{proof}

This transport acts on the two uniform modes and their spectral
labels. Individual incidences and fluctuations within
each sector remain in their complete spaces. The link preserves
exactly the collective datum that the areal weighting will use.

\subsection{Angular character and area operator}
\label{v:ant:sec:ii-area-caracter}

Let $N_e=\operatorname{diag}(0,1,2)$ on each link, and let
$N_{90},N_{120}$ be the sector sums. With the scale
\[
\theta_{\rm clk}=\frac{C^*}{6}\frac{\pi}{180},\qquad
a_0=\frac{1+2\cos(10^\circ)}3\theta_{\rm clk},
\]
fixed, the angular factor is the normalized real character of the triplet
\begin{equation}
 \chi_{10}=\frac13\Tr\operatorname{diag}
       (1,e^{i\pi/18},e^{-i\pi/18})
 =\frac{1+2\cos(10^\circ)}3,
 \qquad a_0=\chi_{10}\theta_{\rm clk}.
 \label{v:ant:eq:ii-area-caracter-triplete}
\end{equation}
The sum of the conjugate phases proves the equality. On the positive branch
considered, \(\theta_{\rm clk}>0\), \(\chi_{10}>0\), and
\(\gamma_{\HMT}>0\); hence \(a_0\) and the areal scale are
positive. The angular triplet and its normalization precede the
entropic distribution on the links. With these quantities,
the normalized area operator is
\[
\widehat{\mathcal A}/\ell_P^2
=\gamma_{\HMT}a_0(N_{90}+N_{120}).
\]
The radial composition fixes \(\ell_P^2=L_\alpha^2\) here, so that
\(\widehat{\mathcal A}=\gamma_{\HMT}a_0L_\alpha^2(N_{90}+N_{120})\).
The same constructed length is transmitted to the area quantum, preserving
the occupation operators and their sectorial recovery.
The spectrum of the normalized operator \(\widehat{\mathcal A}/L_\alpha^2\)
has values $n\gamma_{\HMT}a_0$, $0\le n\le72$,
with multiplicities given by $(1+z+z^2)^{36}$.
The proof consists in summing the eigenvalues $0,1,2$ of the
36 tensor factors.

\begin{theorem}[Recovery of the boundary sectors]
Let
\[
N=N_{90}+N_{120},\qquad D=90N_{90}+120N_{120}.
\]
Then
\begin{equation}
N_{90}=4N-\frac D{30},\qquad
N_{120}=\frac D{30}-3N.
\label{v:ant:eq:recover}
\end{equation}
Thus the total and the weighted reader jointly recover
the two occupations.
\end{theorem}
\begin{proof}
The matrix $\left(\begin{smallmatrix}1&1\\90&120\end{smallmatrix}\right)$
has determinant $30$. Its inversion gives \eqref{v:ant:eq:recover}.
On the domain of integer occupations, the image satisfies
$D\in30\Z$ and $90N\le D\le120N$.
\end{proof}

The modular Hamiltonian in the source is
$K_A=-\log\rho_A=cI+\Delta_{\rm mod}D$.
For $\Delta_{\rm mod}\ne0$ and known normalization,
$D=(K_A-cI)/\Delta_{\rm mod}$ expresses the second observable.
This operator describes the spectral structure of the reduced state;
it is not identified by notation with the generator of temporal evolution.
Area preserves total occupation, and the pair adds provenance.
For example, $(N_{90},N_{120})=(1,0)$ and $(0,1)$ share $N=1$,
but have $D=90$ and $D=120$.

% RESULTADO_RECUPERADO from II REV09. See MAPA_PROCEDENCIA.json.
% Selected lines and explicit mechanical changes; no new proof.
\subsection{State normalization and the meaning of the modular Hamiltonian}
The second observable is made explicit through the state that defines it.
The preceding tensor decomposition is retained: eighteen qutrits
per sector, each with number operator $\operatorname{diag}(0,1,2)$;
$N_m$ is the sum of these operators, and the sectors act on distinct
factors. Let
\[
 z_m=1+e^{-m\Delta_{\rm mod}}+e^{-2m\Delta_{\rm mod}},
 \qquad m\in\{90,120\}.
\]
The state collection in the source has the product realization
\[
 \rho_A=Z^{-1}e^{-\Delta_{\rm mod}(90N_{90}+120N_{120})},
 \qquad Z=z_{90}^{18}z_{120}^{18}.
\]
Commutativity of the number operators and tensor factorization
prove the expression for $Z$. The state has full rank for real,
finite $\Delta_{\rm mod}$, and its spectral logarithm gives
\[
 K_A=-\log\rho_A=(\log Z)I+\Delta_{\rm mod}D.
\]
For $\Delta_{\rm mod}\ne0$, $K_A$ weights the excitations of the
two sectors differently. The joint reading with $N$ recovers the two occupations
through equation~\eqref{v:ant:eq:recover}.

The need for the joint reading is seen in two complementary examples.
The occupation states $(1,0)$ and $(0,1)$ have the same total $N=1$
and different modular weights. Conversely, $(4,0)$ and $(0,3)$ have
the same $D=360$ and different totals. The pair $(N,D)$ separates both
cases. This recovery concerns sectorial occupations, not
each microscopic vector within a degeneracy of those occupations.

When $\Delta_{\rm mod}=0$, the state is proportional to the identity
and $K_A$ is scalar: that reading no longer distinguishes the sectors.
For $\Delta_{\rm mod}\ne0$, the known normalization makes it possible to recover
$D$ and compose it with area. This defines the role of the modular
Hamiltonian: it characterizes the reduced state and its information distribution.
Its inclusion in the article is motivated by that mathematical role,
whereas the evolution Hamiltonian belongs to the dynamical development
of action.


\subsection{Reduced state, Schmidt decomposition, and entanglement entropy}
If $\rho_A=\sum_\nu p_\nu|\nu\rangle\langle\nu|$,
the purification
\[
|\Psi\rangle=\sum_\nu\sqrt{p_\nu}\,
|\nu\rangle_A|\nu\rangle_B
\]
satisfies $\Tr_B|\Psi\rangle\langle\Psi|=\rho_A$.
Thus $-\Tr(\rho_A\log\rho_A)$ is the entanglement entropy
of this particular pure bipartition. The enlarged state declares
which system carries that interpretation.

For a two-sheet involution with one eigenvector of each sign,
\[
\rho_y=\frac{e^{-yR}}{2\cosh y},\qquad R^2=I,
\]
diagonalization gives
\[
S(\rho_y)=\log(2\cosh y)-y\tanh y,\qquad
D(\rho_y\Vert\rho_{-y})=2y\tanh y.
\]
Entropy is even in $y$; the distance explicitly compares
the two orientations. This two-level entropy is not identified
with the capacity $81\log3$ of the triadic cut.
